\section{Literature Review}
\label{sec:related}
The introduction raised two questions, one about the search design for the WFLOP and one about the evidence behind comparisons of such designs. This section reviews the work behind both in five strands. It moves from metaheuristics for the WFLOP (Section~\ref{sec:related-wflop}) to the parameterization of PSO (Section~\ref{sec:related-pso}) and hybrid swarm--local-search designs (Section~\ref{sec:related-hybrid}), and then to benchmarking practice (Section~\ref{sec:related-bench}) and the critiques of metaphor-based algorithms (Section~\ref{sec:related-metaphor}). Each subsection states what is new here, and the last one draws the strands together.

\subsection{Wind Farm Layout Optimization with Metaheuristics}
\label{sec:related-wflop}
\emph{Discrete and continuous formulations.} Early studies placed turbines in the cells of a square grid and optimized the cell occupation with genetic algorithms (GAs)~\cite{Mosetti1994,Grady2005}, and binary PSO with time-varying acceleration coefficients was applied to the same kind of grid~\cite{Pookpunt2013}. A grid makes the problem combinatorial but restricts the layouts that can be found. Kusiak and Song~\cite{Kusiak2010} formulated a continuous model with turbine coordinates in a circular farm, a discretized wind rose, a Jensen wake~\cite{Jensen1983,Katic1986} and explicit boundary and spacing constraints. This model, with two wind data sets~\cite{Kusiak2010,Bansal2017} and several farm radii, has become a common reference for metaheuristics and defines the benchmark used here (Section~\ref{sec:setup}). On it, or on close variants, ant colony optimization (ACO) and particle filtering~\cite{Eroglu2012,Eroglu2013}, biogeography-based optimization (BBO)~\cite{Bansal2017} and random search inside the feasible space~\cite{Feng2015} have been proposed. GA, PSO, differential evolution (DE) and simulated annealing were compared on offshore layouts~\cite{Elkinton2008}, and PSO has reduced wake losses in~\cite{Asaah2021}.

\emph{Newer metaheuristics and hybrids.} Later work used self-informed GAs~\cite{Ju2019}, an adaptive evolutionary algorithm combined with Monte Carlo tree search~\cite{Bai2022}, deterministic refinement techniques~\cite{Nagpal2021} and VNS for large offshore farms~\cite{Cazzaro2022}. Metaphor-based algorithms have also been applied. For example, the water cycle algorithm was compared with the SSA and the grey wolf optimizer on turbine placement~\cite{Rezk2019}, and the Laplacian SSA, previously proposed by Solanki and Deep~\cite{Solanki2023}, was combined with VNS in a preliminary, unpublished comparison of ours that motivated this study. Reviews list many such methods, often evaluated on the same classical wind scenarios but with differing budgets, settings and constraint handling~\cite{Azlan2021}.

\emph{Common benchmarks.} To make results comparable, common benchmarks have been set up. The GECCO layout competitions provided a common evaluation interface, and the top entries of the 2015 edition were compared in~\cite{Wilson2018}. The IEA Wind Task~37 case studies fixed a simplified Gaussian wake, site boundaries and wind roses~\cite{Baker2019,IEA37repo}, and eight methods, gradient-based, gradient-free and hybrid, each run by researchers experienced with it, reached different layouts of similar energy yield on a common problem~\cite{Thomas2023}. Both efforts compare methods \emph{as submitted}: each team chose its own settings and effort, so the results show what experts achieve with their methods, not how the methods compare under one budget and one protocol.

\emph{Gradient-based and model-based alternatives.} Metaheuristics are not the only option. Gradient-based optimization with exact gradients reduces the cost of layout optimization considerably compared with finite differences~\cite{Guirguis2016}. Pseudo-gradients offer a cheap alternative to exact gradients~\cite{Quaeghebeur2021}, and reduced parameterizations of the layout cut the number of design variables~\cite{Stanley2019}. Gaussian wake models~\cite{Bastankhah2014,Tao2020} make the objective smoother than the top-hat Jensen wake. On such smooth models, the cost of a gradient decides how far these methods scale: a forward-difference gradient of the $2N$ coordinates costs $2N+1$ objective evaluations, whereas an analytic or reverse-mode algorithmically differentiated gradient costs a small multiple of one evaluation, at the price of memory~\cite{Rodrigues2024}. With algorithmic differentiation, parallel evaluation of the flow cases and a heuristic initial layout, gradient-based optimization has been applied to farms of 500 turbines~\cite{Rodrigues2024}. Stochastic gradient descent reduces the cost of each step by sampling the wind conditions~\cite{Quick2023}, and a heuristic start followed by gradient-based refinement handles inclusion and exclusion zones~\cite{CriadoRisco2024}. Surrogate models extend layout optimization to larger farms and higher-fidelity flow models~\cite{Yang2023,Wang2024,Li2025}. Graph neural networks, which represent a farm as a graph of interacting turbines and thus apply to any layout, have been trained to predict turbine powers~\cite{ParkPark2019} and wake flow fields~\cite{LiZhang2023}, and have served as surrogates in layout~\cite{ParkPark2019} and yaw optimization~\cite{LiRobert2024}. Finally, with large-eddy simulations as the flow model, Bayesian optimization with Gaussian-process surrogates has optimized layouts and wake steering without gradients~\cite{Bempedelis2024}.

\emph{Position of this study.} We use the Kusiak--Song benchmark in its continuous form (discretized only in the wind directions) and three further wind farm types: a Horns Rev~1 block, a Lillgrund block and IEA37 Case Study~1 (Sections~\ref{sec:farm-I}--\ref{sec:farm-IV}). Unlike competitions, all nine methods (PSO-VNS, PSO, SSA-VNS, SSA, LX-SSA, DE, GA, VNS and MS-SLSQP) are run by the same code under one budget in objective evaluations, which for the gradient-based MS-SLSQP includes its finite-difference gradients. Surrogate-based methods are outside the present pool, because they need training data and are usually reported under their own budgets. Exact gradients need a smooth model~\cite{Guirguis2016,Thomas2023}, whereas the top-hat objective used here is discontinuous. SLSQP with the exact gradient is therefore tested only on the piecewise-smooth IEA37 model, with each gradient charged at its measured cost in evaluations (Section~\ref{sec:iea-grad}). The protocol is nevertheless complementary to such methods. Budgets that count every model evaluation (including gradient and training evaluations), seed-paired runs and equal-cost controls apply to them in the same way; for instance, a heuristic initial layout~\cite{Rodrigues2024,CriadoRisco2024} could be compared with random sampling of equal cost in the farm, as the first phases are here. Among the gradient-free methods of this pool, PSO is the basis of PSO-VNS, and we turn to it next.

\subsection{PSO Parameterization}
\label{sec:related-pso}
Since PSO-VNS uses PSO without modification in its first phase, the parameter setting of PSO needs a justification. The inertia weight~\cite{Shi1998} and the constriction factor of Clerc and Kennedy~\cite{Clerc2002} were introduced to control the divergence of the particles in the PSO of Kennedy and Eberhart~\cite{Kennedy1995}. The constriction setting is equivalent to an inertia-weight PSO with $w\approx0.73$ and $c_1=c_2\approx1.50$. Stability analyses of the particle dynamics~\cite{Poli2009,Cleghorn2018,Bonyadi2017} place this setting inside the region in which the mean and the variance of a particle's position converge (Section~\ref{sec:pso}). We therefore use the standard PSO with the constriction coefficients of Clerc and Kennedy, without tuning, both for PSO and for the first phase of PSO-VNS. With this setting the swarm contracts, and, as noted in Section~\ref{sec:intro}, a contracted swarm refines individual turbines only by chance, which motivates pairing it with a local search.

\subsection{Hybrid Swarm--Local-Search Designs}
\label{sec:related-hybrid}
Combining a population-based global search with a local search is the idea of memetic algorithms. Krasnogor and Smith~\cite{Krasnogor2005} gave a taxonomy and discussed design issues such as the choice of the local search and its integration into the evolutionary cycle, and Neri and Cotta~\cite{Neri2012} reviewed memetic computing for discrete, continuous and constrained problems. VNS~\cite{Mladenovic1997,Hansen2001} changes neighborhoods systematically around an incumbent and has a continuous form~\cite{Mladenovic2008}. In the WFLOP, it has been used for large offshore farms~\cite{Cazzaro2022}, and deterministic refinement techniques for layouts have been compared in~\cite{Nagpal2021}. Outside wind energy, PSO and VNS have already been combined, with VNS as the local search of a PSO, for flowshop scheduling~\cite{Tasgetiren2007}, constrained shortest paths~\cite{Marinakis2017} and sensor-network localization~\cite{Gumaida2019}. The simplest of these designs is the sequential (relay) hybrid, which runs the global method for part of the budget and passes its best solution to the local method.

\emph{Position of this study.} PSO-VNS is a new two-phase variant of this kind of hybrid for the WFLOP. Unlike the PSO--VNS designs above~\cite{Tasgetiren2007,Marinakis2017,Gumaida2019}, which apply VNS within the PSO loop, it runs the two methods in sequence, with exactly counted budgets and a first phase that reproduces a PSO run (Section~\ref{sec:psovns-method}). Its evaluation includes an ablation in which the first phase is compared with equal-cost random sampling in the farm and in its bounding square, and with the local search alone, so that the contribution of each phase is measured rather than assumed. In addition, the second phase is instrumented to show what the local search actually does at the given budget (Section~\ref{sec:ablation}). Measuring such contributions requires a controlled comparison, which is the subject of benchmarking methodology.
% CHECK-FINAL [C53]-[C56]: component analysis verdicts (Section ablation) and CHECK-DIAG [D08]-[D10], [D20] (VNS phase essentially one compass-search descent): stated in Table rq-matrix and the contributions, see the comments there

\subsection{Benchmarking Methodology for Metaheuristics}
\label{sec:related-bench}
Hooker~\cite{Hooker1995} distinguished competitive testing, which tells which algorithm wins, from controlled experiments, which tell why. Later guidelines turned this distinction into practice: equal budgets counted in function evaluations, reported and justified settings of all methods, well-chosen reference algorithms and component analyses of new proposals~\cite{BartzBeielstein2020,LaTorre2021}. Platforms such as COCO~\cite{Hansen2021} standardize budgets, instances and performance measures for continuous black-box optimizers, but on artificial test functions. A constrained engineering problem such as the WFLOP must, in addition, decide how infeasible final solutions enter the comparison.

For the statistical analysis, Dem\v{s}ar~\cite{Demsar2006} and Derrac et~al.~\cite{Derrac2011} recommended nonparametric tests over multiple problems: the Friedman test with post hoc comparisons of average ranks and pairwise Wilcoxon signed-rank tests, with a multiplicity correction such as Holm's~\cite{Holm1979}. Benavoli et~al.\ showed that post hoc tests based on mean ranks depend on the other methods in the pool~\cite{Benavoli2016}, which is why our primary tests of differences use case means; they also advocated Bayesian analyses with a region of practical equivalence~\cite{Benavoli2017}. Carrasco et~al.~\cite{Carrasco2020} reviewed these classical and newer trends for swarm and evolutionary algorithms and gave practical guidelines, and Campelo and Takahashi~\cite{Campelo2019} derived the numbers of instances and runs needed for a given power and accuracy. Finally, equivalence testing with two one-sided tests (TOST) asks whether a difference lies within a margin of practical irrelevance~\cite{Lakens2017}, and thus avoids reading a nonsignificant difference as evidence that two algorithms perform equally.

\emph{Position of this study.} We apply this practice to the WFLOP and combine three elements that are each established elsewhere but that we have not found applied together in WFLOP comparisons. The first is a feasibility-aware ranking of penalty-based methods (qualification, then conditional quality), related to the feasibility rules of constrained evolutionary computation~\cite{Deb2000} and the feasibility-first ranking of the CEC constrained competitions~\cite{Wu2017}. The second is equal-cost random-sampling controls for the first phase of hybrids. The third is equivalence analyses (TOST and a Bayesian signed-rank test) at three levels of inference: over seeds, over cases and over the six designed scenario clusters (Section~\ref{sec:setup}). We did not size the experiment by a power analysis as in~\cite{Campelo2019}; instead, for each equivalence statement, we report the smallest margin at which it holds, so that a reader can apply a different margin.

\subsection{Critiques of Metaphor-Based Algorithms}
\label{sec:related-metaphor}
The weak baselines and uncontrolled comparisons that these guidelines address are also at the center of the critiques of metaphor-based algorithms, which concern the salp-swarm methods in our pool. S\"orensen~\cite{Sorensen2015} argued that many ``novel'' metaheuristics hide known search operators behind a new metaphor and are justified by comparisons with weak baselines, and a group of editors and researchers later called for action on the proliferation of such methods~\cite{Aranha2022}. Camacho-Villal\'on et~al.~\cite{CamachoVillalon2023} showed for six popular metaphor-based algorithms that their search rules reuse known ideas under new names and are presented in a misleading way. For the SSA~\cite{Mirjalili2017}, the critical review of Castelli et~al.~\cite{Castelli2022} analyzed the update rules and the evidence for its performance. In the form used here, the leaders perturb the best solution uniformly with a shrinking radius, and the followers only average (Section~\ref{sec:ssa}).

\emph{Position of this study.} These critiques rest mostly on the update equations or on artificial test functions. We add a component-level test on an applied constrained problem, as part of the ablation of PSO-VNS: whether a salp-swarm first phase does better than random sampling of equal cost in the farm (Section~\ref{sec:ablation}). This test is also a self-critique, as LX-SSA was previously proposed by Solanki and Deep~\cite{Solanki2023}, the first of whom is an author of this study. It assesses LX-SSA as published, including its cost of $2N_p$ evaluations per iteration, and an ablation separates its Laplace step from that cost (Section~\ref{sec:ablation}).

\subsection{Summary and Position of This Study}
Taken together, the five strands show established methods and established protocol elements but, to our knowledge, no WFLOP study that combines a sequential PSO--VNS design with the controls and analyses needed to measure its contribution. Table~\ref{tab:rq-matrix} summarizes the resulting research questions. For each, it gives the closest prior work, the control or analysis missing there and the finding established here. All questions are studied on the problem model stated in Section~\ref{sec:model}.

\begin{table}[!t]
\centering
\caption{Research questions, closest prior work and what this study adds. PSO-VNS (first row) is new; the other methods and the individual protocol elements are established, and the last column gives the evidence established here.}
\label{tab:rq-matrix}
\footnotesize
\begin{tabular}{@{}>{\raggedright\arraybackslash}p{0.18\textwidth}>{\raggedright\arraybackslash}p{0.20\textwidth}>{\raggedright\arraybackslash}p{0.23\textwidth}>{\raggedright\arraybackslash}p{0.31\textwidth}@{}}
\toprule
Research question & Closest prior work & Control or analysis missing there & Finding established here \\
\midrule
Does a sequential PSO--VNS hybrid improve on PSO for the WFLOP? & PSO--VNS hybrids with VNS inside the PSO loop~\cite{Tasgetiren2007,Marinakis2017,Gumaida2019}; PSO on WFLOP~\cite{Pookpunt2013,Asaah2021} & A sequential design with exactly counted phase budgets, compared with PSO on the same seeds & PSO-VNS ranks first on the benchmark; within the margin of PSO with 15$^\circ$ bins, 0.19~pp better with 1$^\circ$ bins (eight cases); feasible more often than PSO on both site blocks (Sections~\ref{sec:farm-I}--\ref{sec:farm-III}) \\
Does the first phase of a hybrid beat random sampling of equal cost? & PSO--VNS and memetic hybrids~\cite{Tasgetiren2007,Marinakis2017,Gumaida2019}; random search in the feasible space~\cite{Feng2015}; SSA critiques~\cite{Castelli2022,CamachoVillalon2023} & Equal-cost random-sampling controls in the site and in its bounding square & A PSO phase gains; salp-swarm phases give no gain larger than 0.05~pp over disc sampling at $4D$ and $5D$ (not at $6D$); the square control flatters them (Sections~\ref{sec:ablation} and~\ref{sec:spacing}) \\
How do unsuccessful runs and constraint handling enter the comparison? & Feasibility rules and CEC rankings~\cite{Deb2000,Wu2017}; WFLOP reviews~\cite{Azlan2021} & Separate reporting of success rates, conditional quality and an all-run outcome; alternative constraint handling & The all-run endpoint keeps the conditional-quality order on the 68 cases (except SSA and MS-SLSQP), not on Lillgrund; in a six-case, five-method study, Deb's rules with box clipping change the order, radial projection keeps the run-level order and the leading method (Sections~\ref{sec:archive-reliability} and~\ref{sec:boundary}) \\
When is a small difference negligible? & Equivalence and Bayesian comparisons~\cite{Lakens2017,Benavoli2016,Benavoli2017,Carrasco2020} & Equivalence at a stated margin, inference level and evaluator, with empirical checks of the error rate & PSO-VNS vs.\ PSO is equivalent on the fixed benchmark (seed level) and under exchangeable-case sensitivity (also in a shifted-null check), not under cluster sensitivity; optimized directly with 1$^\circ$ bins (eight cases), PSO-VNS is better by more than the margin (Sections~\ref{sec:psovns-pso} and~\ref{sec:direction}) \\
Do the rankings hold on other wind farm types and evaluators? & Comparisons as submitted~\cite{Wilson2018,Baker2019,Thomas2023} & The same nine methods, GA included, under one charged budget on four farm types; re-evaluation of layouts & PSO-VNS ranks first on Horns Rev~1 and Lillgrund at 6,030 evaluations (on Lillgrund, MS-SLSQP is better in the paired test) and with 16 IEA37 turbines; GA leads on Horns Rev~1 at 30,030 (also under 1$^\circ$ bins and PyWake) and with 36 IEA37 turbines; on the benchmark, the leader holds with 1$^\circ$ bins (Sections~\ref{sec:farm-I}--\ref{sec:farm-IV}) \\
\bottomrule
\end{tabular}
\end{table}
% CHECK-FINAL [C53]-[C56]: component analysis verdicts (row 2); CHECK-DIAG [D08]-[D10], [D20]: VNS phase essentially one compass-search descent (contributions, item 3)
% CHECK (manual, 2026-10-04): row 3: all-run endpoint on the 68 cases, rank order unchanged except SSA/MS-SLSQP (analysis/rev3_inference.json, tab:S-sa-allrun); constraint study: Deb + clipping order changes, projection unchanged (rev3_constraint.json, tab:S-sa-constraint); row 4: TOST null-calibrated p 0.015 (tab:S-sa-tostaudit), direct 1-deg PSO-VNS - PSO -0.194 pp, case 90% CI [-0.272, -0.114] (rev3_fine.json, tab:S-sa-fine); row 5: HR pool GA best at 30,030, +0.48 (1 deg) / +0.43 GWh/yr (PyWake) (rev3_sites.json, tab:S-sa-hr-pool, tab:S-sa-hr-eval); IEA37 pool: GA first among gradient-free with 36 turbines (tab:S-sa-iea-pool)
